\scanpage{57}
% archive foliation: 56
\deleted{
de façon naturelle un revêtement étale
\deleted{\unsure{fini de $T$ sur $S$}, \unsure{noté}}
\deleted{$\uline{A}(T)$ \unsure{ou} $\bar T$} $f_*^{\mathrm{ét}}(T)$, qui dépend \uline{fonctoriellement} de $T$, \deleted{et transforme} le foncteur $\unsure{f_*^{\mathrm{ét}}|f_!^{\mathrm{ét}}}$ transforme d'ailleurs sommes en \deleted{produits} sommes.

On va définir un $S$-morphisme
\diagram{\begin{tikzcd}T\arrow[r]&\bar T\end{tikzcd}}
i.e. un \unsure{élt|élémt} de $\operatorname{Hom}_S(T,\bar T)$, i.e. une section de $\uline{\operatorname{Hom}}_S(T,\unsure{T|\bar T}$
}

Plus généralement, pour $T$ donné, $S'\mapsto\uline{\operatorname{Hom}}_S(T,S')$ est un foncteur des variétés étales de $S$ dans les rev. étales de $S$, \uline{exact à gauche}.
\deleted{Si \unsure{l'on désirait un} pt géométrique $a$ de $S$, supposé connexe, On obtient ainsi un foncteur}
\inserted{Supposons $S$ connexe.}
\deleted{Un théorème général nous prouve que le foncteur est de la forme $\varinjlim\operatorname{Hom}($}
\margin{(\unsure{inutile?}) \unsure{raisonner directement} sur le foncteur $S'\mapsto\operatorname{Hom}_S(T,S')$]}
Fixons un pt géométrique $a$ de $S$ supposé connexe (Ce qui ne devrait pas être nécessaire) \deleted{\unsure{Alors la}} la définition
